% SPDX-License-Identifier: Apache-2.0
\section{A Neighbour's Register}
\label{sec:x-sorted}

A sorted list kept in registers moves its bytes along by one place
when a byte is put in: each place after the new byte takes the byte of
the place before it.
In a loop over the places that is \code{self.l[j - 1]}, the loop's
variable moved by a constant, which the lowering folds into a register's
name as it unrolls the loop: \code{l\_1} reads \code{l\_0}, and so on
(issue 1851).
A loop's variable plus or minus a constant is the only arithmetic an
index takes, as in Section~\ref{sec:x-regs} it is a number or the
loop's variable; an index that the unrolled loop would make below zero
is refused, naming the index and the variable's value.

\code{Sorted<N>} keeps the \code{N} largest bytes it is given, largest
first.
A byte in goes after every byte not below it; the places before keep
theirs, its place takes it, and the places after take their
neighbours', the last byte falling off.
Its second loop starts at 1, since place 0 has no place before it.
The netlist of \code{Sorted<4>} is checked against the run under nvc
and Verilator, as Section~\ref{sec:x-checked} describes.

\begin{figure}[htbp]
\centering
\input{sorted_timing.tex}
\caption{Five bytes in: each goes to its place and the bytes after it
move along by one; the fifth, larger than the rest, pushes the
smallest off the end.}
\label{fig:x-sorted}
\end{figure}

\lstinputlisting[caption={\code{ex\_sorted.rs}. Run with \code{bazel run //lib/examples:ex\_sorted}.},label={lst:x-sorted}]{lib/examples/ex_sorted.rs}

\lstinputlisting[style=ascii,caption={What \code{ex\_sorted} printed when the build ran it: the places after each byte, the first, then the lowering.},label={lst:x-sorted-out}]{out_sorted.txt}
